\chapter{Error Handling and Robustness}
\label{ch:error-handling}

\textit{Expecting the unexpected without crashing the system.}

C++ provides two orthogonal mechanisms: \textbf{assertions} for catching programmer bugs at development time, and \textbf{exceptions} for handling runtime environmental failures gracefully. This chapter covers both, along with RAII's role in exception safety, the standard exception hierarchy, custom exception types, and debugging strategies for C++98/03.

%% ============================================================
\section{The Philosophy of Failure}
\label{sec:philosophy-failure}

\begin{center}
\begin{tabular}{|l|l|l|}
\hline
\textbf{Failure type} & \textbf{Cause} & \textbf{Mechanism} \\
\hline
Bug               & Programmer violated a precondition & \texttt{assert} — crash in debug, removed in release \\
Runtime error     & Environmental failure               & Exception — survive and recover \\
\hline
\end{tabular}
\end{center}

Never use exceptions as a control-flow mechanism for normal program paths.

%% ============================================================
\section{Assertions}
\label{sec:assertions}

\texttt{assert} from \texttt{<cassert>} evaluates a condition at runtime. If false, it prints a diagnostic and calls \texttt{abort()}. In release builds (\texttt{-DNDEBUG}), all \texttt{assert} calls compile out to zero overhead.

\begin{lstlisting}[language=C++, caption={Using assert to enforce preconditions}]
#include <iostream>
#include <cassert>

int divide(int a, int b) {
    assert(b != 0 && "Denominator cannot be zero!");
    return a / b;
}

int main() {
    std::cout << divide(10, 2) << "\n";
    return 0;
}
\end{lstlisting}

\begin{lstlisting}[language=C++, caption={Conditional debug logging macro}]
#include <iostream>

#ifdef DEBUG
    #define LOG(x) std::cout << "DEBUG: " << x << "\n"
#else
    #define LOG(x)
#endif

int main() {
    int x = 10;
    LOG("x is " << x); // Only printed with -DDEBUG
    return 0;
}
\end{lstlisting}

%% ============================================================
\section{Basic \texttt{try-catch-throw}}
\label{sec:try-catch-throw}

\begin{lstlisting}[language=C++, caption={Basic exception flow}]
#include <iostream>
#include <stdexcept>

void connect_to_server() {
    bool network_down = true;
    if (network_down)
        throw std::runtime_error("No internet connection.");
    std::cout << "Connected!\n";
}

int main() {
    try {
        connect_to_server();
    }
    catch (const std::runtime_error& e) {
        std::cerr << "Error: " << e.what() << "\n";
    }
    return 0;
}
\end{lstlisting}

%% ============================================================
\section{Catching Multiple Exception Types}
\label{sec:multiple-catch}

\begin{lstlisting}[language=C++, caption={Multiple catch clauses — specific before general}]
#include <iostream>
#include <stdexcept>

void risky(int choice) {
    if (choice == 1) throw std::out_of_range("Index out of bounds");
    if (choice == 2) throw std::runtime_error("Network failure");
    if (choice == 3) throw 42;
}

int main() {
    for (int i = 1; i <= 4; ++i) {
        try { risky(i); }
        catch (const std::out_of_range& e) {
            std::cerr << "Range error: " << e.what() << "\n"; }
        catch (const std::runtime_error& e) {
            std::cerr << "Runtime error: " << e.what() << "\n"; }
        catch (int e) {
            std::cerr << "Integer exception: " << e << "\n"; }
        catch (...) {
            std::cerr << "Unknown exception!\n"; }
    }
    return 0;
}
\end{lstlisting}

%% ============================================================
\section{The Standard Exception Hierarchy}
\label{sec:std-exceptions}

\begin{lstlisting}[language={}, caption={Standard exception hierarchy}]
std::exception (base)
├── std::logic_error
│   ├── std::invalid_argument
│   ├── std::domain_error
│   ├── std::length_error
│   └── std::out_of_range
└── std::runtime_error
    ├── std::range_error
    ├── std::overflow_error
    └── std::underflow_error
std::bad_alloc
std::bad_cast
std::bad_typeid
\end{lstlisting}

\begin{lstlisting}[language=C++, caption={Catching via the base class}]
#include <iostream>
#include <stdexcept>

int main() {
    try {
        throw std::out_of_range("Index 99 out of [0,10)");
    }
    catch (const std::exception& e) { // Catches any standard exception
        std::cerr << "Caught: " << e.what() << "\n";
    }
    return 0;
}
\end{lstlisting}

%% ============================================================
\section{Custom Exception Classes}
\label{sec:custom-exceptions}

\begin{lstlisting}[language=C++, caption={Simple custom exception}]
#include <exception>
#include <iostream>

class DatabaseError : public std::exception {
public:
    virtual const char* what() const throw() {
        return "Database connection failed";
    }
};

int main() {
    try { throw DatabaseError(); }
    catch (const std::exception& e) { std::cerr << e.what() << "\n"; }
    return 0;
}
\end{lstlisting}

\begin{lstlisting}[language=C++, caption={Custom exception with structured error data}]
#include <stdexcept>
#include <string>
#include <iostream>

class AppError : public std::runtime_error {
    int error_code;
    int line_number;
public:
    AppError(const std::string& msg, int code, int line)
        : std::runtime_error(msg), error_code(code), line_number(line) {}
    virtual ~AppError() throw() {}
    int getCode()       const throw() { return error_code; }
    int getLineNumber() const throw() { return line_number; }
};

int main() {
    try { throw AppError("Parse failed", -12, 42); }
    catch (const AppError& e) {
        std::cerr << e.what()
                  << " (code=" << e.getCode()
                  << ", line=" << e.getLineNumber() << ")\n";
    }
    return 0;
}
\end{lstlisting}

%% ============================================================
\section{Exception Specifications (C++98)}
\label{sec:exception-specifications}

\begin{lstlisting}[language=C++, caption={C++98 exception specifications}]
void func()     throw(int)  { throw 42; } // May throw int
void no_throw() throw()     {}            // No exceptions (deprecated in C++11)
void anything()             {}            // May throw anything
\end{lstlisting}

If a function with \texttt{throw(T)} throws an unlisted type, \texttt{std::unexpected()} is called, which by default calls \texttt{std::terminate()}. Exception specifications were \textbf{deprecated in C++11} and \textbf{removed in C++17}; replaced by \texttt{noexcept}.

%% ============================================================
\section{Stack Unwinding and RAII}
\label{sec:stack-unwinding}

When an exception is thrown, C++ exits each scope between the \texttt{throw} and the matching \texttt{catch}, calling the \textbf{destructor} of every local object in reverse order.

\begin{lstlisting}[language=C++, caption={Stack unwinding guarantees destructors run}]
#include <iostream>
#include <stdexcept>

class Resource {
    const char* name;
public:
    explicit Resource(const char* n) : name(n) {
        std::cout << name << ": Acquired\n";
    }
    ~Resource() { std::cout << name << ": Released\n"; }
};

void risky() {
    Resource r("File handle");
    throw std::runtime_error("Boom!");
}  // r's destructor IS called during unwinding

int main() {
    try { risky(); }
    catch (...) { std::cout << "Caught exception\n"; }
    return 0;
}
/* Output:
   File handle: Acquired
   File handle: Released
   Caught exception */
\end{lstlisting}

\textbf{Raw \texttt{new}/\texttt{delete} inside a function is dangerous} --- the \texttt{delete} is skipped on exception. Wrap resources in RAII classes or \texttt{std::vector}.

%% ============================================================
\section{Rethrowing Exceptions}
\label{sec:rethrow}

\begin{lstlisting}[language=C++, caption={Correct rethrow preserves dynamic type}]
#include <iostream>
#include <stdexcept>

void log_and_rethrow() {
    try {
        throw std::runtime_error("Original error");
    }
    catch (const std::exception& e) {
        std::cerr << "Logging: " << e.what() << "\n";
        throw;     // Correct: rethrows original type
        // throw e; // WRONG: slices to std::exception copy!
    }
}

int main() {
    try { log_and_rethrow(); }
    catch (const std::exception& e) {
        std::cerr << "Final: " << e.what() << "\n";
    }
    return 0;
}
\end{lstlisting}

%% ============================================================
\section{Best Practices: Throw by Value, Catch by \texttt{const\&}}
\label{sec:throw-catch-rules}

\begin{lstlisting}[language=C++, caption={Correct throw/catch pattern}]
try {
    throw std::runtime_error("Disk full"); // Throw by VALUE
}
catch (const std::runtime_error& e) {     // Catch by CONST REFERENCE
    std::cout << e.what() << "\n";
}
\end{lstlisting}

\begin{itemize}
    \item \textbf{No copy}: catching by reference avoids constructing a copy of the exception object.
    \item \textbf{No slicing}: catching a derived type by reference preserves all derived-class data via polymorphism.
    \item \textbf{No pointer ownership}: never throw \texttt{new ExceptionType()} --- the catcher cannot safely \texttt{delete} it.
\end{itemize}

%% ============================================================
\section{Nested Exceptions and Function Try Blocks}
\label{sec:nested-exceptions}

\begin{lstlisting}[language=C++, caption={Nested try-catch}]
#include <iostream>
#include <stdexcept>

int main() {
    try {
        try {
            throw std::out_of_range("inner error");
        }
        catch (const std::out_of_range& e) {
            std::cerr << "Inner: " << e.what() << "\n";
            throw; // Propagate
        }
    }
    catch (const std::exception& e) {
        std::cerr << "Outer: " << e.what() << "\n";
    }
    return 0;
}
\end{lstlisting}

\begin{lstlisting}[language=C++, caption={Constructor function try block}]
#include <stdexcept>
#include <iostream>

class Widget {
    int* data;
public:
    Widget(int n)
    try : data(new int[n])
    {
        if (n == 0) throw std::invalid_argument("Size must be > 0");
    }
    catch (const std::bad_alloc& e) {
        std::cerr << "Allocation failed: " << e.what() << "\n";
        throw;
    }
    ~Widget() { delete[] data; }
};

int main() {
    try { Widget w(0); }
    catch (const std::exception& e) {
        std::cerr << "Caught: " << e.what() << "\n";
    }
    return 0;
}
\end{lstlisting}

%% ============================================================
\section{Debugging Strategies for C++98/03}
\label{sec:debugging}

\subsection{GDB / LLDB Command Reference}

\begin{center}
\begin{tabular}{|l|l|}
\hline
\textbf{Command} & \textbf{Effect} \\
\hline
\texttt{break file.cpp:42}  & Set breakpoint at line 42 \\
\texttt{run}                & Start the program \\
\texttt{print varname}      & Print variable value \\
\texttt{backtrace} / \texttt{bt} & Show call stack \\
\texttt{watch varname}      & Pause when variable changes \\
\texttt{step} / \texttt{next}  & Step into / step over \\
\hline
\end{tabular}
\end{center}

\subsection{Compile-Time Assertions (C++98)}

\begin{lstlisting}[language=C++, caption={C++98 compile-time assertion trick}]
template<bool B> struct CompileTimeCheck {};
template<>       struct CompileTimeCheck<true> { typedef void type; };

// Fails at compile time if sizeof(int) != 4
typedef CompileTimeCheck<sizeof(int) == 4>::type INT_IS_4_BYTES;
\end{lstlisting}

C++11 provides the cleaner \texttt{static\_assert(condition, "message")}.

%% ============================================================
\section{Professional Insights: Exception Safety and Design}
\label{sec:exception-safety}

\subsection{Exception Safety Guarantees}

\begin{center}
\begin{tabular}{|l|l|}
\hline
\textbf{Guarantee} & \textbf{What it means} \\
\hline
\textbf{No-throw}  & Never throws. Marked \texttt{throw()} in C++98 or \texttt{noexcept} in C++11. \\
\textbf{Strong}    & If it throws, no observable state change (rollback semantics). \\
\textbf{Basic}     & If it throws, no resources are leaked; objects are in a valid state. \\
\hline
\end{tabular}
\end{center}

\subsection{\texttt{noexcept} Forward Reference (C++11)}

C++11 replaces \texttt{throw()} with \texttt{noexcept}:

\begin{lstlisting}[language=C++, caption={noexcept marks non-throwing functions}]
void increment(int& x) noexcept { x++; }
// If it does throw: std::terminate() is called immediately.
// Compiler skips all exception-unwinding bookkeeping.
\end{lstlisting}

\subsection{Exception Hierarchy Design Rules}

\begin{itemize}
    \item Derive all application exceptions from \texttt{std::exception} so catch-all handlers use \texttt{const std::exception\&}.
    \item Use \texttt{std::runtime\_error} for recoverable failures; \texttt{std::logic\_error} for programming errors.
    \item \textbf{Never throw from destructors.} A destructor called during stack unwinding that itself throws results in \texttt{std::terminate()}.
\end{itemize}

\subsection{When Not to Use Exceptions}

Exceptions carry overhead for setting up unwinding tables and branch-testing on every entry to a try block. They are inappropriate for:
\begin{itemize}
    \item Functions on hot paths called millions of times per second (HFT, game physics).
    \item Signal handlers and interrupt service routines.
    \item Deeply embedded systems where unwinding tables inflate binary size.
\end{itemize}

In these contexts, prefer error codes (\texttt{int}, \texttt{enum}) or sentinel return values with explicit caller-side checking.
